Can a single set of weights act as a ``computer''?  
We term this abstraction a Neural Computer (NC): a neural network that models computer interface dynamics learned from experience.
% An NC couples environment dynamics, working memory, and interface I/O through pixels, actions, and language.
% In this paper, we train neural networks (NNs) to learn the dynamics of traditional computers (TCs). The NN's inputs include screenshots, instructions (natural language or program commands), and user actions such as keyboard typing and mouse movements/clicks.


To implement this idea, we instantiate NCs as video models.
Several lines of evidence suggest that, at this stage, video models are the most practical substrate for prototyping this vision—though we expect the long-term solution to require a fundamentally new neural architecture (\Cref{section:toward-cnc}).
%
This choice builds on perspectives like world models~\citep{ha2018world}, which demonstrate that neural networks can internalize environment dynamics and support predictive imagination; 
and advances in high-capacity video generators such as Veo~3.1~\citep{google_veo3_1_2025} and Sora~2~\citep{openai_sora2_2025} show that such learned dynamics can be rendered into coherent, photorealistic frame sequences; besides, 
frontier interactive video models such as Genie~3~\citep{bruce2024genie} further extend this trajectory, demonstrating action-controllable generative environments.
% 
In parallel, LLM-driven UI systems such as Imagine with Claude\footnote{\url{https://claude.ai/imagine/}} map natural-language inputs to structured interface updates. 
%` 
Yet these capabilities remain split across components: models~render~or~predict, while an external execution environment (simulator/OS/DOM) carries executable state and mediates control.

%To address this gap, we define a neural computer (NC) as a learned, executable, and ultimately programmable system in which the model, environment, and runtime share a unified latent state. An NC exposes a computer-like interface through pixels, actions, and language. 
We study two interfaces that instantiate our NC formulation (see ~\Cref{section:prelim}). 
\nccligen{} is a terminal interface conditioned on text (natural language or command lines) and an initial frame (\Cref{section:impl-cligen}). 
\ncguiworld{} is a desktop interface conditioned on recent pixels and synchronized mouse/keyboard actions (\Cref{section:impl-guiworld}).  
% 

\begin{tcolorbox}[
  colback=metablue!8,
  colframe=metablue!40,
  boxrule=0pt,
  borderline west={1pt}{0pt}{gray!60},
  left=6pt,right=4pt,top=3pt,bottom=3pt]
\small
\textbf{Neural Computer (NC) abstraction~(\hyperlink{teaser}{Teaser}).}
A neural system $(F, G)$ parameterized by $\theta$ that models the dynamics of an interactive computer interface 
% (iii) renders and controls the interface via pixels and actions
with a single latent runtime state $h_t$ that acts as its working memory (see~\cref{eq:wm}).
\end{tcolorbox}


In the \nccligen{} experiments, the NC can render and execute basic command-line workflows.
It often stays aligned with the terminal buffer and captures common ``physics'' of everyday CLI use (e.g., fast scrollback, prompt wrapping, window resizing).
On carefully scripted data, rollouts can be visually and structurally close to real sessions.
The model can execute short command chains and render their outputs.
Arithmetic-probe scores improve substantially with stronger system-level conditioning.
%Reprompting boosts probe accuracy from 4\% to 83\% (\Cref{fig:cligen-exp6}).

In the \ncguiworld{} experiments, we evaluate standard world-model designs across action injection, action encoding, and data quality.
Figure~\ref{fig:nc-general} summarizes this design template across two interface-specific NCs trained separately without shared parameters.
Qualitatively, the model learns coherent pointer dynamics and short-horizon action responses (e.g., hover/click feedback and window/menu transitions).
%With explicit cursor supervision, it reaches 98.7\% cursor accuracy (\Cref{tab:cursor-loss}).

Our experimental insights indicate that current NCs already realize several primitives of conventional computers, most notably I/O alignment and short-horizon control.
We define the long-term goal, a Completely Neural Computer (CNC), as a fully learned computer whose compute, memory, and interfaces are unified in a single learned runtime substrate rather than engineered as separate modules.
However, these prototypes represent only an early step toward this CNC vision.
Substantial challenges remain in achieving robust long-horizon reasoning, reliable symbolic processing, and stable interface rendering.
\Cref{section:toward-cnc} presents a CNC capability map and outlines key open challenges toward the CNC vision. 

\begin{tcolorbox}[
  colback=metablue!8,
  colframe=metablue!40,
  boxrule=0pt,
  borderline west={1pt}{0pt}{gray!60},
  left=6pt,right=4pt,top=3pt,bottom=3pt]
\small
\textbf{Completely Neural Computer (CNC) abstraction~(\Cref{sec:roadmap}).}
A Neural Computer instance is \emph{complete} (i.e., a CNC) if it is
(i) \emph{Turing complete},
(ii) \emph{universally programmable},
(iii) \emph{behavior-consistent} unless explicitly reprogrammed, and
(iv) realizes the architectural and programming-language semantics of NCs relative to conventional computers.
\end{tcolorbox}

\textbf{Concretely, this work makes the following contributions:}
{
\renewcommand\labelitemi{\large$\bullet$}
\begin{itemize}
\setlength{\itemsep}{2pt plus 1pt minus 1pt}
\setlength{\parskip}{0pt}
\item Define neural computers (NCs) and build video-based prototypes for both CLI and GUI interfaces.
% \item Release open-source video datasets spanning thousands of hours across both domains.
\item Provide a data engine and alignment recipe that synchronize text, actions, and frames~for~CLI and~GUI~environments.
\item Identify practical design choices for NCs through extensive ablation studies.
\item Outline a roadmap toward completely neural computers (CNCs) grounded in primitives of conventional computers and open research questions.

\end{itemize}
}
